\section{Recovering coordinates using points alone}
\label{sec:frame-inverse}

Coordinates can be recovered from betweenness and segment congruence. This
gives a converse to the interpretation of geometry by scalar pairs and
specifies the exact sense of equivalence. Throughout this section
$B(A,M,D)$ is inclusive betweenness and $AB\equiv CD$ means segment
congruence. All formulas defining geometric constructions below use just
these relations, equality, and three named points. A line or a scalar
never has to be added as a primitive geometric object.

\subsection{A unit frame and four construction formulas}

Write $L(A,D,P)$ if one of $B(A,D,P),B(D,P,A),B(P,A,D)$ holds. Write
$M(A,H,D)$ if $B(A,H,D)$ and $AH\equiv HD$. Define
\[
 V(A,D;C,E)\quad\Longleftrightarrow\quad
 \exists H\,[M(A,H,E)\land M(D,H,C)].
\]
In a coordinate plane over a linearly ordered field $F$, $M(A,H,D)$
means $2H=A+D$. Indeed, for unequal endpoints put
$H=A+t(D-A)$ with $0\leq t\leq1$. Equality of squared lengths gives
$t^2=(1-t)^2$ after cancellation of $\|D-A\|^2>0$, hence $2t=1$.
For equal endpoints, betweenness forces $H=A=D$. The midpoint always
exists since $2$ is invertible. Thus
\[
 V(A,D;C,E)\quad\Longleftrightarrow\quad D-A=E-C.
\]
This single predicate expresses translation even when all four points
lie on one line or a displacement vanishes.

For distinct $C,D$ define a direction test
\[
 T(P,Q;C,D)\quad\Longleftrightarrow\quad
 C\ne D\ \land\ \exists E\,[V(C,D;P,E)\land L(P,E,Q)] .
\]
Here $E=P+(D-C)\ne P$, so the test means that $Q-P$ is some scalar
multiple of $D-C$. It includes $Q=P$. For completeness, on an affine
line any three points are ordered by their one-dimensional parameters.
Their middle parameter gives one of the three displayed betweenness
relations. For $E\ne P$ a parameter for $Q$ can be computed by dividing
along a nonzero coordinate of $E-P$. This establishes both directions
of the claim using decidable order.

Choose points $O,I,J$ with $O\ne I$, $OI\equiv OJ$, and with a point
$I'$ such that $M(I,O,I')$ and $JI\equiv JI'$. Call this formula
$\mathrm{Frame}(O,I,J)$. Set $u=I-O$, $v=J-O$, $N=u\cdot u$.
The midpoint equation gives $I'=O-u$. The last congruence gives
$4u\cdot v=0$, so $u\cdot v=0$ and $v\cdot v=N>0$. The identity
\[
 \det(u,v)^2+(u\cdot v)^2=(u\cdot u)(v\cdot v)
\]
gives $\det(u,v)\ne0$. Thus $u,v$ form a basis. The frame
$(0,0),(1,0),(0,1)$ witnesses existence in every such plane.
Reflection of a frame is allowed.

The scalar domain is the line $H=\{x:L(O,I,x)\}$, with vertical line
$K=\{U:L(O,J,U)\}$. Lowercase variables $x,y,z$ below are restricted
to $H$. The following formulas define transfer, addition, multiplication,
and the coordinates of a point.
\begin{align*}
 S(x,U)&\ \Longleftrightarrow\ U\in K\land T(x,U;I,J),\\
 A(x,y,z)&\ \Longleftrightarrow\ V(O,x;y,z),\\
 P(x,y,z)&\ \Longleftrightarrow\
   \exists U,W\,[S(y,U)\land S(z,W)\land T(x,W;I,U)],\\
 G(Q;x,y)&\ \Longleftrightarrow\
   \exists U\,[S(y,U)\land V(O,x;U,Q)].
\end{align*}
The additive identity is $O$ and the multiplicative identity is $I$. More explicitly, if $x=O+a u$, then
$S(x,U)$ determines $U=O+a v$. Indeed, writing $U=O+bv$, the vector
$-a u+bv$ is a multiple of $v-u$ precisely when $a=b$.
If $y=O+b u$, then $A$ gives $z=O+(a+b)u$. For multiplication, put
$W=O+c v$ and $U=O+bv$. The test
$T(x,W;I,U)$ says that $-a u+c v$ is proportional to $-u+bv$.
Equality of the first coefficients forces the proportionality factor
to be $a$. The second coefficients give $c=ab$. The fixed horizontal
component $-1$ in the reference vector ensures $I\ne U$ for
every $b$. This proves the product formula also for zero and negative
inputs. Finally $G(Q;x,y)$ says exactly $Q=O+a u+b v$.

\begin{figure}[htbp]
\centering
\begin{tikzpicture}[x=1.15cm,y=.84cm,font=\small,
  axis/.style={->,gray!70}, drawline/.style={line width=.85pt,color=accent},
  pt/.style={circle,fill=ink,inner sep=1.55pt}]
\begin{scope}
  \draw[axis] (-.2,0)--(2.75,0);\draw[axis] (0,-.22)--(0,2.30);
  \draw[drawline,dashed] (0,1)--(1,0);
  \draw[drawline] (0,1.65)--(1.65,0);
  \node[pt,label=below left:$O$] at (0,0) {};
  \node[pt,label=below:$I$] at (1,0) {};
  \node[pt,label=left:$J$] at (0,1) {};
  \node[pt,label=below:{$(b,0)$}] at (1.65,0) {};
  \node[pt,label=left:{$(0,b)$}] at (0,1.65) {};
  \node[anchor=north,font=\sffamily\small] at (1.02,-.62) {Transfer to the second axis};
\end{scope}
\begin{scope}[xshift=6.35cm]
  \draw[axis] (-.2,0)--(2.50,0);\draw[axis] (0,-.22)--(0,3.05);
  \draw[drawline,dashed] (0,1.65)--(1,0);
  \draw[drawline] (0,2.64)--(1.6,0);
  \node[pt,label=below left:$O$] at (0,0) {};
  \node[pt,label=below:$I$] at (1,0) {};
  \node[pt,label=below:{$(a,0)$}] at (1.6,0) {};
  \node[pt,label=left:{$(0,b)$}] at (0,1.65) {};
  \node[pt,label=left:{$(0,ab)$}] at (0,2.64) {};
  \node[anchor=north,font=\sffamily\small] at (1.00,-.62) {Read the product};
\end{scope}
\end{tikzpicture}
\caption{Multiplication in the recovered field. In each diagram the solid line
is parallel to the dashed line. The transfer construction moves the final
ordinate back to the baseline. The formulas also cover zero and negative
inputs.}\label{fig:intercept}
\end{figure}


Define $R(z)$ by $B(O,z,I)\lor B(O,I,z)$, and define
$x\preceq y$ by $\exists z\,[A(x,z,y)\land R(z)]$. For
$z=O+c u$, the first disjunct gives $0\le c\le1$, the second gives
$1\le c$, and their union gives $c\ge0$. This recovers the directed
order on the whole baseline. Additive inverse and division by a nonzero
scalar have unique graphs from addition and multiplication. When a total
reciprocal is part of the scalar language, assign reciprocal zero at zero.

\subsection{The two inverse maps}

\paragraph{Framed interpretation theorem.}
For every decidably linearly ordered field $F$, constructing its
point plane and interpreting the baseline recovers an isomorphic ordered
field. Constructing the point plane of that baseline recovers a point
structure isomorphic to the original framed plane. Both isomorphisms
have explicit graphs. For a primitive recursive presentation of $F$
with decidable equality, order and primitive recursive totalized
division, the two isomorphisms on codes are primitive recursive.

\emph{Proof.}
The scalar isomorphism is $a\mapsto O+a u$ by the four computations
above. For the point isomorphism use $G$. Its inverse sends a point
$Q$ to the two baseline points indexed by
\[
 a=\frac{(Q-O)\cdot u}{N},\qquad
 b=\frac{(Q-O)\cdot v}{N}.
\]
Indeed, the difference $Q-O-au-bv$ is orthogonal to both basis vectors
and hence zero by the nonzero determinant. For any two such points their
squared distance in the original plane equals
$N((a_1-a_2)^2+(b_1-b_2)^2)$. Since the factor $N>0$ is common to all
segments, congruence is preserved and reflected. Affine interpolation
commutes with the coordinate map and its inverse, preserving and
reflecting inclusive betweenness. The maps carry the three named points
to the canonical frame. The displayed expressions use a fixed finite
number of field operations and one division by a proved positive
element. On a coded field each step respects code equality. This also
supplies the asserted primitive recursive implementations.\hfill$\square$

For example, testing which side of a directed line a point lies on
uses the sign of a determinant after a frame has been chosen. Reversing
the orientation of the frame reverses this sign. Equidistance and
betweenness survive every frame change, including a common change in the
length unit. Naming the frame specifies the scalar identification used
in the theorem.

\subsection{An intrinsic point-only axiomatization}

There is a direct way to state the scope of the inverse theorem for
abstract structures. Let $\mathcal F$ be a first-order theory of
decidably linearly ordered fields. For example it may require square
roots of all positive elements, or in addition roots of every odd-degree
polynomial. The latter choice gives real-closed ordered fields. Form a
point theory with $B,\equiv,O,I,J$ by imposing the following explicit
axioms. Every mention of $H,K,S,A,P,G$ abbreviates the point formulas
above.
\begin{enumerate}
\item The named triple satisfies the frame formula. On $H$, the graphs
$A,P$ are everywhere defined and single-valued. With constants
$O,I$, operations so defined, and relation $\preceq$, the relativized
axioms of $\mathcal F$ hold.
\item For every $Q$ there is exactly one pair $(x,y)\in H^2$ with
$G(Q;x,y)$, and for every such pair there is exactly one $Q$.
Here quantification over a pair means two ordinary point quantifiers.
\item For points with coordinates supplied by $G$, congruence is equality
of the two squared distances in $H$. For coordinate pairs $a,b,c$,
$B(a,b,c)$ is equivalent to
\[
 \det(b-a,c-a)=0\quad\land\quad(b-a)\cdot(b-c)\le0 .
\]
Both are fixed polynomial tests over the interpreted field, hence
first-order point formulas using the graphs of arithmetic.
\end{enumerate}
To check the betweenness test, if $a=c$ the inequality forces $b=a$
by nonnegativity of squares. Otherwise the determinant supplies
$b=a+t(c-a)$ by division along a nonzero component. The inequality is
$t(t-1)\|c-a\|^2\le0$, equivalently $0\le t\le1$. Thus all three
axioms hold in the plane of every field satisfying $\mathcal F$.

Conversely, in any point structure satisfying this list the interpreted
baseline is a model of $\mathcal F$ by item 1. Item 2 supplies the graph
of an isomorphism onto its coordinate plane, and item 3 verifies both
directions on each primitive predicate. To see that the named frame is preserved,
item 3 first makes the coordinate map an isomorphism of unframed point
structures, and therefore carries the original defining formulas to the
same formulas in an actual field plane. There the calculations above
give $G(O;O,O)$, $G(I;I,O)$ and $G(J;O,I)$ (with the frame and scalar
points carried along by that map). Pulling these formulas back gives
the assertion in the original structure. This proves a two-sided
first-order interpretation for the stated presentation. Uniquely
defined operations may equivalently be retained as graphs throughout,
so the proof requires no choice of representatives or choice principle.
For the point language without named constants, the same formulas
give an interpretation with a frame as parameter.

The list is a coordinatization axiom presentation in point language.
Its algebraic and geometric content is explicit. A representation theorem
for a separately chosen synthetic axiom system with intuitionistic logic
would require deriving the list from that system. The inverse theorem here
applies directly to the stated presentation and to the exact plane used
in this report.

\paragraph{Checked Lean component.}
\texttt{FrameRelations.lean} defines the point-language graphs.
\texttt{CanonicalFrame.lean} proves the transfer, addition, multiplication,
order and coordinate formulas in the standard frame over every linearly
ordered field, and proves that the coordinate graph is a bijection.
Its executable algebraic test for $B(A,M,D)$ uses collinearity together
with $|AM|^2\le |AD|^2$ and $|MD|^2\le |AD|^2$.
For unequal endpoints the parameter $t$ converts these inequalities to
$t^2\le1$ and $(1-t)^2\le1$. Their conjunction is precisely
$0\le t\le1$. Equal endpoints force $M=A$. The file verifies that
equivalence, so its point predicates agree with those above. The Lean
results concern the field semantics. The interpretation of the underlying
arithmetic in iPRA is a separate part of the argument.
